\documentclass[11pt]{article}
\usepackage[left=1.25in,right=1.25in,top=1.2in,bottom=1.2in]{geometry}
\usepackage{lmodern}\usepackage[T1]{fontenc}\usepackage{microtype}
\usepackage[hyphens]{url}
\setlength{\parskip}{7pt}\setlength{\parindent}{0pt}
\pagestyle{empty}
\begin{document}

\hfill 7 October 2026

\medskip
Dear Editor,

We resubmit ``Does AI rate through the cycle? Procyclicality in credit assessment'' to
\emph{Finance Research Letters}. The three reports were detailed and have improved the paper; all
sixteen comments are answered point by point. The revision adds a real-fundamentals arm using
FY2024 data for 80 SEC filers, where the procyclical slope replicates in every model with all
intervals excluding zero; two frontier models, taking the design to seven across two families; and
three further replicate seeds re-measured 2.5 months later, bounding coefficient drift at $0.037$
notches. The central finding is unchanged: every model is significantly procyclical, and an
explicit ``through-the-cycle'' instruction anchors the rating while suppressing the default
probability that should keep moving.

Where a referee's request produced a result that does not help us, we report it---issuer recall
identifies seven of forty real firms, which we decline to read as equivalence, and the three
extra-seed captures refused by the pre-registered error-rate gate are disclosed with their causes.
A complete replication archive accompanies the paper: one command regenerates every reported
number and binds it to the artifact that produced it
(\url{doi.org/10.5281/zenodo.21864041}). A tracked-changes copy is included.

Thank you for considering the revision.

Sincerely,

Samir Chincholikar \\
Robin Chawla

\medskip
Corresponding author: \\
Robin Chawla \\
\texttt{robin.chawla.cse14@iitbhu.ac.in}

\end{document}
